\section{Empirical eigenvector laws and conditional entropy}\label{sec:entropy}

Fix \(E\in I_{d,\kappa}\), and abbreviate
\(\rho=\rho_{d,w}(E)>0\).  The local law describes the environment seen
by a typical vertex.  We now retain eigenvector coordinates at that
vertex as well.  The objective of this section is an entropy inequality
for their law conditional on the entire potential environment.  In
Section~\ref{sec:rigidity}, equality in a complementary entropy bound
will identify this conditional law as Gaussian.

Choose a measurable real orthogonal eigenbasis \(u_1,\ldots,u_n\),
ordered by eigenvalue and normalized by
\(\langle u_i,u_j\rangle=n\delta_{ij}\).  Write
\begin{equation}\label{eq:microscopic-positions}
 x_i=n(\lambda_i-E),\qquad \nu_n=\sum_{i=1}^n\delta_{x_i}.
\end{equation}
A finite list of eigenvectors gives each vertex an \(\RR^m\)-valued
mark, one coordinate for each member of the list.  We call these
coordinates colors.  All lists below are finite, but their length is
arbitrary and imposes no additional smallness condition on \(w\).

\subsection{Keeping the empirical law in the limit}

For a fixed finite list, sample a uniform vertex and read its marked
neighborhood.  On a tree neighborhood use a uniformly ordered rooted
tree isomorphism; outside it use the nonbacktracking lift, recording a
failure flag.  Conditional on the graph, potentials, and eigenvectors,
this procedure defines an empirical probability law.  It is important
to retain this law, rather than just its average: independent vertex
samples are independent conditional on the empirical law.

\begin{proposition}[Marked subsequential limits]\label{prop:marked-limits}
Every sequence of admissible sizes tending to infinity has a
subsequence on which \(\nu_n\) converges in law to a locally finite
point measure \(\nu\).  On a further subsequence one can retain jointly
the positions and the empirical marked-neighborhood laws for the
eigenvalue-ordered lists in an exhaustion by bounded intervals,
together with all their finite truncations and restrictions.
Repeated limiting positions retain separate marks.

Almost every resulting empirical law \(\mathcal L\), with \(m\) colors,
is an automorphism-invariant law on the \(d\)-regular tree, whose
potential marginal is the deterministic iid uniform law on
\([-w,w]\).  If \(U_v\in\RR^m\) is its mark at \(v\), then
\begin{equation}\label{eq:mark-normalization-equation}
 \EE_{\mathcal L}U_oU_o^{\mathsf T}=\Id_m,
 \qquad
 (V_v-E)U_v+\sum_{u\sim v}U_u=0\quad\text{at every }v.
\end{equation}
The empirical laws and their restrictions to smaller lists can be
chosen consistently.
\end{proposition}

\begin{proof}
The Wegner bound of Lemma~\ref{lem:spectral-averaging} gives tightness
of the point measures on each compact interval.  Choose interval
endpoints uncharged almost surely by a subsequential limit; such an
exhaustion exists, since the expected limiting counts are bounded by
the same constant times interval length.  Truncate each list length
and then stratify by that length.  At a uniformly sampled vertex the
sum of squared coordinates of a list of length \(m\) has mean \(m\).
This proves tightness of all finite-radius observations.  A diagonal
extraction retains all radii, list truncations, and interval
restrictions, as well as their positions.

For each fixed radius the fraction of vertices with a non-tree
neighborhood tends to zero.  Uniform rooted orderings give invariance
under root-preserving automorphisms.  Sampling an oriented edge and
reversing it gives invariance under re-rooting across an edge; together
these give full tree-automorphism invariance.  The empirical law of
the potential labels alone converges in probability to their iid
tree law: one and two sampled bounded neighborhoods have this law
asymptotically, and two disjoint neighborhoods use independent
potentials.  Thus this marginal is deterministic inside each
limiting \(\mathcal L\).

The second-moment assertion needs more than tightness.  If the list
comes from \([-L,L]\), spectral decomposition gives
\begin{equation}\label{eq:marks-dominated-by-resolvent}
 \sum_{i\text{ in list}}u_i(v)^2
 \le (1+L^2)\Im G_{vv}(E+i/n).
\end{equation}
The right side is uniformly integrable at a sampled vertex by
Corollary~\ref{cor:local-consequences}.  Consequently the limiting
empirical second moments lose no mass.  More explicitly, represent
the joint convergence almost surely, first pass bounded
second-moment truncations, and then remove the truncation using the
tail bound in expectation.  The empirical moments before taking the
limit are exactly \(\delta_{ij}\), so their limits are the same almost
surely; polarization handles the mixed moments.  The eigenvalue
equations pass to the limit on every finite neighborhood because
\(\lambda_i=E+O_L(n^{-1})\).  Countable consistency gives the assertions
simultaneously at all vertices and for all retained lists.
\end{proof}

We henceforth fix one such \(\mathcal L\) when making assertions about
a conditional law.  We write \(\omega=(V_v)_{v\in T_d}\) for its full
potential environment; thus \(\omega\) here records the coupled
potentials themselves.  Expectations with respect to its deterministic
iid marginal are denoted \(\EE_\omega\).

\subsection{Eliminating a conditional mean}

Orthogonality alone does not show that
\(\EE_{\mathcal L}[U_o\mid\omega]\) vanishes.  A coherent component
depending on the potential environment would survive it.  The
following perturbation averages such a component over a physical
energy interval longer than the microscopic window, while changing
the ensemble by a vanishing total variation distance.

\begin{lemma}[Conditional centering]\label{lem:conditional-mean}
For almost every empirical law in Proposition~\ref{prop:marked-limits},
\begin{equation}\label{eq:conditional-zero-mean}
 \EE_{\mathcal L}[U_v\mid\omega]=0
 \qquad\text{at every vertex }v.
\end{equation}
\end{lemma}

\begin{proof}
Let \(f\) be a bounded Lipschitz function of the potentials in a
fixed rooted ball.  Average over its branch orderings and put the
corresponding finite-graph value \(f(v)\) equal to zero at tree
failures.  First require that \(f\) vanish unless the central
potential belongs to a compact subinterval of \((-w,w)\).  Choose
\(b_0\in C_c^\infty((-w,w))\), nonnegative and bounded below by a
positive constant on that subinterval.  For
\(0\le t\le a_n:=n^{-2/3}\), change every potential by
\[
 V_v\longmapsto V_v+t b_0(V_v).
\]
For large \(n\) this is a diffeomorphism of \([-w,w]\), identical
near its endpoints.  The one-site density ratio is \(1+O(t)\), has
integral one, and has second moment \(1+O(t^2)\).  Independence
therefore bounds the total variation distance of the product laws by
\(O(\sqrt n\,t)=o(1)\), uniformly over this range.  The normalized
test vector \(f/\sqrt n\) changes by \(O_f(t)\) in Euclidean norm.

For the original configuration set
\(B=\diag(b_0(V_v))\).  The support condition writes
\(f/\sqrt n=B^{1/2}q_f\), with \(\|q_f\|\le C_f\).  Scalar spectral
averaging gives, for every interval \(I\),
\begin{equation}\label{eq:bump-spectral-averaging}
 \int_0^{a_n}
 \left\langle f/\sqrt n,\mathbf1_I(H+tB)f/\sqrt n\right\rangle\,dt
 \le C_f|I|.
\end{equation}
For completeness, compress the resolvent by Schur complement to the
indices where \(B>0\), and conjugate by \(B^{1/2}\).  It becomes
\((t\Id-T)^{-1}\) with \(\Im T>0\) at positive spectral height.
The integral of its imaginary part in a vector, over the real
\(t\)-axis, is at most \(\pi\) times that vector's squared norm.
The Herglotz representation, followed by spectral inversion, proves
\eqref{eq:bump-spectral-averaging}; restricting the nonnegative
integrand to \([0,a_n]\) preserves the bound.

Take \(I=E+[-L,L]/n\), divide by \(a_n\), and average over the
original data.  Total variation comparison and the change in the
test vector show that
\[
 \EE\left\langle f/\sqrt n,\mathbf1_I(H)f/\sqrt n\right\rangle
 \le O_f(n^{-1/3})+o(1).
\]
The left side is
\(\EE\sum_{x_i\in[-L,L]}(n^{-1}\sum_v f(v)u_i(v))^2\).
Hence every empirical root average against \(f\) vanishes in the
limit.  Passing these first moments is justified by the uniform
second-moment bounds.  A countable dense collection of tests on all
finite balls, followed by an exhaustion of \((-w,w)\) in the central
coordinate, determines the full conditional mean.  Averaging tests
over branch orderings loses no information, since the rooted joint
law is invariant under those permutations.  Re-rooting proves the
claim at every vertex.
\end{proof}

\subsection{Counting labels while retaining the environment}

We use natural logarithms, write \(H\) for discrete entropy, and
write \(\mathcal H\) for differential entropy.  Conditional
differential entropy means the integral of the entropy of the
conditional law.  Our entropy argument follows the colored-tree and
Gaussian-smoothing framework of Backhausz and Szegedy
\cite[Theorems~6.1 and 8.1]{BackhauszSzegedy2019}.  The extra issue is
that the conditioning environments attached to adjacent vertices
overlap.  We keep these labels in the count, and later compute their
entropy contribution explicitly.

Fix a radius \(R\ge0\) and divide \([-w,w]\) into \(Q\) equal
intervals.  The environment label at a vertex records the rooted
isomorphism type of its radius-\(R\) neighborhood with these binned
potentials, including its central bin.  Adjoin a label from any
fixed finite alphabet for its mark.  Tree failures may use a separate
label, which still records the central potential bin.  The resulting
alphabet has fixed finite size \(M\).
For any vertex labeling let \(T\) be its ordered star histogram,
symmetrized over neighbor permutations, and let \(S\) be its
oriented-edge histogram.

With probability tending to one over the graph and potentials,
simultaneously for all such histograms, the number of labelings
whose prescribed central bins are correct is at most
\begin{equation}\label{eq:environment-type-count}
 \exp\left\{n\left[H(T)-\frac d2H(S)-\log Q+o(1)\right]\right\}.
\end{equation}
The requirement that an entire environment label be correct can only
decrease this number.  Here is the first-moment count, including the
normalization of the potential term.  In the configuration model,
prescribe each vertex label and the labels desired at its ordered
half-edges.  If \(m_{ab}=ndS(a,b)\), the probability of a compatible
matching is
\[
 \frac{\prod_{a<b}m_{ab}!\,\prod_a(m_{aa}-1)!!}{(nd-1)!!},
\]
and is zero unless \(m_{ab}=m_{ba}\) and all \(m_{aa}\) are even.
Its logarithm is \(-ndH(S)/2+O_M(\log n)\).  The arrays of each
ordered star type \(T_0\) number at most \(\exp(nH(T_0))\), with
\(H(T_0)\le H(T)\).  There are only polynomially many types for
fixed \(M\).  Independently, the prescribed central bins agree with
the iid potentials with probability exactly \(Q^{-n}\).
Stirling's formula, Markov's inequality with a sublinear exponential
allowance, and a union bound prove \eqref{eq:environment-type-count}.
Conditioning on simplicity changes the estimates by a bounded factor.
This is the fixed-alphabet form of the count in
\cite[Section~3.1, Lemma~3.2]{bu}; its argument is independent of the
eigenvectors and remains valid with the added labels.

\subsection{From the count to full conditional entropy}

Fix \(\eta>0\) before taking the graph-size limit.  Add to the mark
array a centered Gaussian field, with independent colors and
covariance in each color
\begin{equation}\label{eq:entropy-smoothing}
 \Gamma_\eta=\Im(H-E-i\eta)^{-1}
     =\frac{\eta}{(H-E)^2+\eta^2}.
\end{equation}
Let \(F^\eta\) be the limiting sum of marks and noise.  At this fixed
height polynomial approximation on the uniformly bounded spectra
identifies the local noise covariance with the infinite-tree
covariance \(\Im(H_{T_d}-E-i\eta)^{-1}\).  Given \(\omega\), the
noise is independent of the marks.  Fix neighbor indices
\(1,\ldots,d\) on the rooted tree used for \(\mathcal L\).  A star
observation consists of the root followed by these neighbors; an
edge observation consists of the root and neighbor \(1\).  Their
ordering is inherited from the uniformly ordered finite lift,
and is the same ordering used in \(\omega\).

\begin{proposition}[Entropy conditional on the tree environment]
\label{prop:conditional-entropy}
For almost every empirical law \(\mathcal L\) with \(m\ge1\) colors,
and for every fixed height in any prescribed countable set of
positive heights,
\begin{equation}\label{eq:conditional-star-edge-entropy}
 \begin{split}
 &\mathcal H(F^\eta_{\rm star}\mid\omega)
       -\frac d2\mathcal H(F^\eta_{\rm edge}\mid\omega)\\
 &\hspace{1cm}\ge
 m\left[\frac12\log(2\pi\mathrm e\,\eta)
       -\int\log|t-E-i\eta|\,d\nu_{d,w}(t)\right].
 \end{split}
\end{equation}
\end{proposition}

\begin{proof}
For fixed environment radius \(R\) and number of bins \(Q\), let
\(A_{R,Q}\) and \(B_{R,Q}\) be the tuples of environment labels
seen by the ordered star and edge, respectively.  The count first gives
an inequality with these finite conditioning data:
\begin{align}
 &m\left[\frac12\log(2\pi\mathrm e\,\eta)
        -\int\log|t-E-i\eta|\,d\nu_{d,w}(t)\right]\notag\\
 &\quad\le
 \mathcal H(F^\eta_{\rm star}\mid A_{R,Q})
 -\frac d2\mathcal H(F^\eta_{\rm edge}\mid B_{R,Q})
 +H(A_{R,Q})-\frac d2H(B_{R,Q})-\log Q.
 \label{eq:finite-environment-entropy}
\end{align}
We shall derive this inequality by removing the mark bins at fixed
\(R,Q\).  The remaining environment contribution is
\begin{equation}\label{eq:environment-entropy-remainder}
 H(A_{R,Q})-\frac d2H(B_{R,Q})-\log Q.
\end{equation}
Its limit is zero as \(Q\to\infty\) at fixed \(R\); after that
cancellation, increasing \(R\) gives conditioning on \(\omega\).

We first keep the environment radius, its bins, and the noisy-mark
bins fixed.  Conditional on the finite graph, potentials, and list,
only the Gaussian noise is random.  The entropy of its binned
labeling is at most the entropy of its histogram type plus the
logarithm of the number of labelings of that type.  The first term
is \(O(\log n)\).  Using \eqref{eq:environment-type-count} for the
second term therefore bounds the entropy per vertex by
\[
 \EE_{\rm noise}H(T)-\frac d2\EE_{\rm noise}H(S)-\log Q+o(1).
\]
We may replace the random histograms by their conditional means in
this expression.  Indeed their individual cell frequencies have
variance tending to zero.  To see this, \(\Tr\Gamma_\eta^2=O_\eta(n)\)
implies that the average squared cross-covariance between two
independently sampled bounded local blocks is \(o(1)\).  Their
marginal covariances have deterministic upper and positive lower
bounds at fixed \(\eta\).  Gaussian continuity consequently makes
the average total variation distance to the product of the two
marginals tend to zero.  This remains true after arbitrary shifts
by the fixed eigenvectors and after binning.  Overlapping blocks
have vanishing probability.  Thus all histogram entropies converge
to those of the corresponding limiting observations.  Gaussian
smoothing removes mark-bin boundary atoms, and the potential bins
have null boundaries as well.

We next remove the mark binning, while retaining the finite
environment labels.  Use cubes of side \(\Delta\) in \(\RR^m\),
clipped outside a large box.  Before clipping, the entropy of the
global quantization is at least the global Gaussian differential
entropy minus \(nm\log\Delta\): the entropy within any cube is at
most its log volume.  Clipping changes the entropy per vertex by
an amount tending to zero with the box size, at fixed \(\Delta\).
For this last assertion it suffices to bound the average second
moment, which is bounded by normalization of the marks and the
fixed-height Gaussian covariance bound.  For example, compare the
tail-bin distribution with summable lattice weights proportional
to \((1+|j|)^{-C}\); its entropy is bounded by its tail probability
times a logarithmic moment bound.  The same argument applies to
each finite-dimensional limiting observation.

After removing the clipping, send \(\Delta\downarrow0\).  The
smoothed densities are bounded and continuous, and have finite
second moments.  Replacing a density by its uniform filling in
each bin gives \(L^1\) convergence; the density bound and the
preceding tail estimate give entropy convergence.  The log-volume
terms cancel because
\(m(d+1)-(d/2)(2m)=m\).  The lower bound furnished by the full
Gaussian array is
\begin{align*}
 \frac1n\mathcal H(\text{global Gaussian noise})
 &=\frac m2\log(2\pi\mathrm e)
       +\frac m{2n}\log\det\Gamma_\eta\\
 &=m\left[\frac12\log(2\pi\mathrm e\,\eta)
       -\frac1n\sum_{i=1}^n\log|\lambda_i-E-i\eta|\right].
\end{align*}
The empirical spectral measure converges to \(\nu_{d,w}\), by
polynomial approximation and convergence of bounded-neighborhood
observations.  Splitting each limiting joint entropy into the entropy of its
environment labels and its conditional mark entropy gives
\eqref{eq:finite-environment-entropy}.

We now compute the environment contribution
\eqref{eq:environment-entropy-remainder} before increasing the radius.  Set
\[
 s_R=\sum_{k=0}^{R}b^k,\qquad p_R=\sum_{k=0}^{R-1}b^k,
 \qquad p_0=0.
\]
The star's collection of labels reads \(1+ds_R\) potential bins:
the ball of radius \(R+1\), with its first neighbor positions
ordered and its remaining branches unordered.  The edge reads
\(2s_R\) bins in its two outward half-balls.  When all these bins
are distinct the overlaps identify the vertices uniquely.  There
are \(dp_R\) independent outward permutations of \(b\) branches
in the star reading and \(2p_R\) in the edge reading.  Thus
\begin{align*}
 H(A_{R,Q})&=(1+ds_R)\log Q-dp_R\log(b!)+o_Q(1),\\
 H(B_{R,Q})&=2s_R\log Q-2p_R\log(b!)+o_Q(1).
\end{align*}
For fixed \(R\), bin collisions have probability \(O_R(Q^{-1})\);
their entropy cost is \(O_R(Q^{-1}\log Q)\), which justifies the
errors.  Both the bin counts and the permutation losses cancel in
\eqref{eq:environment-entropy-remainder}, leaving \(o_Q(1)\).

Take nested potential partitions as \(Q\to\infty\).  Almost surely
all finitely many potentials in a reading are distinct, so its
limiting information is exactly the corresponding continuous
potential configuration modulo the indicated permutations.
Conditioning instead on fully ordered potentials in that union
does not change the observation entropy: each omitted permutation
fixes the observed star or edge, and invariance of \(\mathcal L\)
makes its conditional observation law identical under that
permutation.  Finally let \(R\to\infty\); these ordered
configurations exhaust \(\omega\).

Here the conditioning limits can be justified without subtracting
infinite entropies of continuous environments.  Use relative
entropy with respect to a fixed product Gaussian reference on
the observation coordinates and apply monotone convergence on
the increasing conditioning sigma-fields.  The conditional
relative entropy is finite: fixed-height smoothing bounds
conditional entropy below by noise entropy, while second moments
bound it above.  This proves convergence of the conditional
entropies and hence \eqref{eq:conditional-star-edge-entropy}.

All passages have been made for fixed finite choices before
refinement.  To make them hold inside almost every empirical law,
include the countably many bin choices and heights in the joint
subsequence construction, and take a further subsequence on which
the high-probability count bounds hold eventually at every fixed
precision.  An almost sure representation of the empirical data
can retain the original finite data by sampling their conditional
law.  Their deterministic relations, including the entropy
counts and eigenvector normalization, are then preserved.  The
preceding limits apply to each represented realization and give
the stated almost sure assertion.
\end{proof}

The right side of \eqref{eq:conditional-star-edge-entropy} diverges
as \(\eta\downarrow0\).  This divergence has a specific source: the
eigenvalue equation confines the unsmoothed star marks to one
hyperplane per color.  The next section removes precisely those
normal directions, computes the remaining Gaussian constants,
and uses the resulting finite entropy inequality to identify the
marks.
